% ===================================================================
%  도표 제 3 호 — 어린 시절과 젊은 시절.
%
%  Traduction, faite en 2026, en coreen d'aujourd'hui, sur l'ido de
%  texto/io. Regles de la colonne : voir 10-dopyo-01.tex. Deux
%  scenes, deux numerotations : les cles portent c1 et c2.
%
%  DIX-SEPT RETOURNEMENTS SUR VINGT-CINQ PAIRES, et le tableau
%  explique pourquoi : c'est un DEFILE, et un defile se decrit par
%  attributs. La nourrice AVEC sa coiffe, le mendiant APPUYE sur sa
%  bequille, la vieille APPUYEE sur son baton, le meunier AVEC son
%  chapeau de feutre, les abeilles DE la ruche, la fleche DU clocher,
%  le clocher QUI couronne l'eglise. Chacun de ces liens est une
%  inversion en coreen, et il y en a dix-sept. Le tableau 2, qui
%  enumerait, en coutait un seul. Meme opposition qu'entre les
%  tableaux 13 et 15 telougous, et pour la meme raison.
%
%  LE FAC-SIMILE DONNE DEUX FOIS LE RENVOI (18) ET DEUX FOIS LE (19)
%  au bloc c1-09-1 : les arbres fruitiers puis les pechers pour le
%  premier, le presbytere puis les abricotiers pour le second. C'est
%  une faute de l'edition ido, et la source ne bouge pas : la
%  traduction porte les quatre, dans l'ordre recu, comme les
%  trente-huit colonnes qui precedent.
%
%  UN MOT QUI EST DEVENU UNE INSULTE, TROISIEME FOIS DU RELEVE. La
%  baraque foraine du bloc c2-14-1 montre « la giganto e la nano ».
%  Le coreen a 난쟁이 pour le second, qui traduit exactement, et qui
%  sert aujourd'hui d'injure courante. La colonne ecrit donc
%  아주 작은 사람. Meme raison que pour le కసాయి telougou du tableau
%  15 : un mot juste en 1926 devenu une injure en 2026 ne se met pas
%  dans une traduction de 2026 sans qu'on l'ait voulu. La baraque,
%  elle, reste — c'est la CHOSE, et on ne la deplace pas.
%
%  LE COREEN A LES QUATRE-VINGTS OBJETS DE CE TABLEAU, et deux sont a
%  noter parce qu'ils vont plus loin qu'un mot : 나막신 (75) sont les
%  sabots de bois, chaussure coreenne reelle et pas un calque ; 요람
%  (26) est le berceau, 벌통 (21) la ruche, 목발 (67) la bequille,
%  은방울꽃 (37) le muguet — « la fleur clochette d'argent », qui dit
%  la chose mieux que le mot francais. Trois emprunts seulement, tous
%  pour des objets venus d'Europe : 포플러 (3), 플라타너스 (17),
%  코르셋 (32).
% ===================================================================

%%K t03-apar-1 apar
\VUsaut{3.0mm}
\VUtitre{15.0pt}{60}{도표 제 3 호}
\VUsaut{1.6mm}
\VUtitre{11.4pt}{0}{어린 시절과 젊은 시절.}
\VUsaut{3.4mm}

%%K t03-apar-2 apar
(제 3 호와 제 4 호 도표는, 네 \VUgras{집안 마당}에서 \VUgras{날씨}, \VUgras{나이},
\VUgras{집안}, \VUgras{옷차림}, \VUgras{형편}에 관한 기본 개념을 보여 주도록 서로
엮여 있다.)

\VUsaut{4.0mm}
%%K t03-tit-1 sub
\VUcentre{12.0pt}{0}{\textit{첫째 마당.}}
\VUsaut{3.0mm}

%%K t03-tit-2 sub
\VUpk{10.2pt}{40}{시골 마을의 세례식.}
\VUsaut{2.4mm}

%%K t03-tit-3 sub
\VUpk{10.2pt}{40}{봄.}
\VUsaut{3.0mm}

%%K t03-c1-01-1 p
1. --- 지금은 \VUgras{봄}, \VUgras{삼월}이나 \VUgras{사월}이나 \VUgras{오월},
\VUgras{주}의 여러 날 가운데 \VUgras{월요일}이나 \VUgras{화요일}이나
\VUgras{수요일}이다. 아침 \VUgras{열 시}다.

%%K t03-c1-02-1 p
2. --- \VUgras{날씨}는 좋고 \VUgras{공기}는 맑다. 해가 빛난다. \VUgras{조각구름}
\textsuperscript{(2)} 몇 점이 떠오른다 — 푸른 \VUgras{하늘} \textsuperscript{(1)}로.

%%K t03-c1-03-1 p
3. --- \VUgras{나무}에는 아직 \VUgras{잎}이 거의 없다. 가느다란 \VUgras{포플러}
\textsuperscript{(3)}와 키 큰 \VUgras{플라타너스} \textsuperscript{(17)}에는 지금까지
움만 있을 뿐이다.

%%K t03-c1-04-1 p
4. --- \VUgras{제비} \textsuperscript{(4)}가 돌아와 즐겁게 지저귀며 \VUgras{첨탑}
\textsuperscript{(7)} 둘레를 날아다닌다. 그 첨탑은 \VUgras{종탑} \textsuperscript{(5)}의
것이고, 그 종탑은 마을 \VUgras{교회} \textsuperscript{(6)} 위에 관처럼 얹혀 있다.

%%K t03-tit-4 sub
\VUsaut{3.4mm}
\VUpk{10.2pt}{40}{교회.}
\VUsaut{3.0mm}

%%K t03-c1-05-1 p
5. --- 그 교회는 아주 수수하다. 큰 \VUgras{정문} \textsuperscript{(13)}과 커다란
\VUgras{시계} \textsuperscript{(8)}가 있다. 짧은 \VUgras{바늘} \textsuperscript{(9)}은
숫자 X 위에 있어 \VUgras{시}를 가리킨다. \VUgras{긴 바늘} \textsuperscript{(10)}은 XII
위에 있고(정오), \VUgras{분}을 가리킨다.

%%K t03-c1-06-1 p
6. --- 종탑에서는 \VUgras{종} \textsuperscript{(11)}이 즐겁게 울린다. 지붕 꼭대기에는
돌로 된 작은 \VUgras{십자가} \textsuperscript{(12)}가 서 있다.

%%K t03-c1-07-1 p
7. --- 교회에는 큰 \VUgras{정문} \textsuperscript{(13)}만 있는 것이 아니라 옆에 작은
문도 있다. 그 작은 문으로 \VUgras{본당 신부} \textsuperscript{(15)}님이
\VUgras{수단}을 입으신 채 \VUgras{제의실} \textsuperscript{(14)}에서 나와
\VUgras{사제관} \textsuperscript{(19)}으로 가신다.

%%K t03-c1-08-1 p
8. --- 그 곁에 \VUgras{우유 파는 아낙} \textsuperscript{(24)}이 있다 — 자기
\VUgras{나귀} \textsuperscript{(23)}와 함께. 그 아낙은 아이들에게 줄 좋은 우유를 가져다
준 도시에서 돌아오는 길이다.

%%K t03-tit-5 sub
\VUsaut{4.0mm}
\VUpk{10.2pt}{40}{과수원.}
\VUsaut{3.4mm}

%%K t03-c1-09-1 p
9. --- 알리보롱 씨(나귀)는 그 아침 걸음이 아주 흡족하다. 그는 \VUgras{꽃}
\textsuperscript{(20)} 향기가 밴 공기를 달게 들이마신다 — 그 꽃들은 이웃
\VUgras{과수원}의 \VUgras{과일나무} \textsuperscript{(18)}를 덮고 있다. 그
\VUgras{복숭아나무} \textsuperscript{(18)}와 \VUgras{살구나무} \textsuperscript{(19)}는
온통 분홍빛과 흰빛이어서 좋은 수확을 바라게 한다.

%%K t03-c1-10-1 p
10. --- 그동안 작은 \VUgras{꿀벌} \textsuperscript{(22)}들은 — \VUgras{벌통}
\textsuperscript{(21)}에서 나온 것들이다 — 윙윙거리며 제 달콤한 \VUgras{꿀}을 거두러
그리로 몰래 간다.

%%K t03-c1-11-1 p
11. --- 그런데 무슨 일인가? 마을 아이들이 달려와 놀란 \VUgras{거위} \textsuperscript{(25)}
들을 쫓는다. 공증인의 아이가 \VUgras{세례}를 받고 나서 행렬이 교회에서 나오는 것이다.

%%K t03-c1-12-1 p
12. --- 어린 야코부스는 제 \VUgras{요람} \textsuperscript{(26)} 안에서 즐거운 종소리에
깨어난다. 그의 누이 막달레나도 세례 행렬을 보고 싶어서, 집 문지방에서 머리를 빗기고 옷을
입혀 달라고 어머니께 조른다.

%%K t03-c1-13-1 p
13. --- 어머니가 그러마고 하신다. 어머니는 \VUgras{머릿수건} \textsuperscript{(28)},
\VUgras{윗도리} \textsuperscript{(29)}, \VUgras{앞치마} \textsuperscript{(30)}를 걸치고
문 앞의 작은 걸상에 나와 앉으신다. 막달레나는 \VUgras{속치마} \textsuperscript{(31)}와
\VUgras{코르셋} \textsuperscript{(32)}만 입고 있다.

%%K t03-c1-14-1 p
14. --- 그 아이가 얼마나 기쁜지! 그 아이는 제가 좋아하는 꽃들을 잊는다 — 제
\VUgras{패랭이꽃} \textsuperscript{(34)}, 그 위에 앉는 \VUgras{나비} \textsuperscript{(35)},
제 \VUgras{튤립} \textsuperscript{(36)}, 제 \VUgras{은방울꽃} \textsuperscript{(37)} —
그것은 \VUgras{화분} \textsuperscript{(38)}에 담겨 \VUgras{담} \textsuperscript{(33)} 위에
있다 —, 그리고 그토록 고운 울타리를 이루는 제 \VUgras{장미} \textsuperscript{(39)}를.
그 아이는 행렬에 선 부인들의 화려한 옷차림을 눈여겨본다.

%%K t03-c1-15-1 p
15. --- 마을 아이들은 옷차림에는 별로 마음 쓰지 않는다. 그들은 \VUgras{사탕}
\textsuperscript{(55)}이나 \VUgras{동전} \textsuperscript{(52)}을 얻으려고 서로 밀치며
달려든다. 그 가운데 하나가 운다 \textsuperscript{(40)}.

%%K t03-c1-16-1 p
16. --- 다른 아이는 \VUgras{개} \textsuperscript{(42)}의 발을 밟았다. 개는 깽깽거리며
달아나면서 \VUgras{암탉} \textsuperscript{(43)}과 그 \VUgras{병아리} \textsuperscript{(44)}
들을 쫓아 버린다.

%%K t03-tit-6 sub
\VUsaut{5.0mm}
\VUpk{10.2pt}{40}{세례 행렬.}
\VUsaut{4.0mm}

%%K t03-c1-17-1 p
17. --- 행렬 앞에는 \VUgras{유모} \textsuperscript{(45)}가 걸어간다. 그 여자는 긴 댕기가
달린 고운 \VUgras{머리쓰개} \textsuperscript{(46)}를 썼다. 그 여자는 \VUgras{갓난아이}
\textsuperscript{(47)}를 안고 있는데, 아이는 온통 \VUgras{레이스} \textsuperscript{(48)}로
감싸여 있다.

%%K t03-c1-18-1 p
18. --- 유모 뒤로는 \VUgras{대부} \textsuperscript{(49)}와 \VUgras{대모}
\textsuperscript{(53)}가 온다. 그들이 사탕과 동전을 나누어 준다. 대부는 \VUgras{장갑}
\textsuperscript{(51)}을 끼고 \VUgras{실크해트} \textsuperscript{(50)}를 썼다. 대모는
손에 고운 \VUgras{꽃다발} \textsuperscript{(54)}을 들었다.

%%K t03-c1-19-1 p
19. --- 그들 뒤로는 아이의 \VUgras{아저씨} \textsuperscript{(56)}와
\VUgras{아주머니} \textsuperscript{(58)}가 보인다. 아주머니는 맵시 있는 \VUgras{모자}
\textsuperscript{(59)}를 썼고, 아저씨는 검은 \VUgras{프록코트} \textsuperscript{(57)}에
흰 조끼 차림이다. 그 다음이 \VUgras{사촌 오빠} \textsuperscript{(60)}와
\VUgras{사촌 언니} \textsuperscript{(61)}다. 사촌 언니는 갓난아이의 \VUgras{누이}
\textsuperscript{(62)}인 어린 베르타의 손을 잡고 있다.

%%K t03-c1-20-1 p
20. --- 베르타의 다른 \VUgras{오라비} \textsuperscript{(63)}는 뒤에서 걷는다. 그는
\VUgras{친척 아주머니} \textsuperscript{(64)}에게 손을 내민다. 집안의 \VUgras{친구들}
\textsuperscript{(65)}은 그 뒤에 온다.

%%K t03-tit-7 sub
\VUsaut{4.6mm}
\VUpk{10.2pt}{40}{구경하는 사람들.}
\VUsaut{3.4mm}

%%K t03-c1-21-1 p
21. --- 행렬 가까이에 \VUgras{거지} \textsuperscript{(66)}가 하나 있다 — 자기
\VUgras{목발} \textsuperscript{(67)}에 몸을 기댄 채. 그는 동냥을 청한다.

%%K t03-c1-22-1 p
22. --- 다른 쪽에는 아주 \VUgras{예의 바른} \textsuperscript{(68)} 사내아이가 있다.
그는 아는 사람들에게 인사하며 「\VUgras{안녕하세요}」 하고 말한다.

%%K t03-c1-23-1 p
23. --- 마을 사람들은 세례 행렬을 보려고 집에서 나왔다. \VUgras{시골 사람}
\textsuperscript{(69)}이 \VUgras{무명 모자}를 쓰고 있고 그 곁에 \VUgras{시골 아낙}
\textsuperscript{(70)}이 있다. 다음으로 \VUgras{할머니} \textsuperscript{(71)}가 있는데
자기 \VUgras{지팡이} \textsuperscript{(72)}에 몸을 기댔다. 끝으로 \VUgras{방앗간 주인}
\textsuperscript{(73)}이 있는데 큰 \VUgras{펠트 모자} \textsuperscript{(74)}를 썼고,
\VUgras{나막신} \textsuperscript{(75)}을 신은 젊은 시골 아낙이 제 어린것에게 걸음마를
가르치고 있다.

%%K t03-c1-24-1 p
24. --- 그 마당은 아주 재미있다.

\VUsaut{4.0mm}
%%K t03-tit-8 sub
\VUcentre{12.0pt}{0}{\textit{둘째 마당.}}
\VUsaut{3.0mm}

%%K t03-tit-9 sub
\VUpk{10.2pt}{40}{마을 잔치.}
\VUsaut{2.4mm}

%%K t03-tit-10 sub
\VUpk{10.2pt}{40}{물놀이와 배 겨루기.}
\VUsaut{3.0mm}

%%K t03-c2-01-1 p
1. --- \VUgras{여름}이 왔다. \VUgras{유월}(칠월이나 \VUgras{팔월})의 뜨거운 해가
젊은이들을 미역 감고 배 젓는 일로 이끈다.

%%K t03-c2-02-1 p
2. --- 강 오른쪽 기슭에서 배 젓는 이들이 물놀이와 배 겨루기에 나가려고 옷을 벗는다.

%%K t03-c2-03-1 p
3. --- 그 가운데 하나가 제 \VUgras{신} \textsuperscript{(1)}을 벗는다. 끈을 풀고
단추를 끄른다. 그의 동무 둘은 벌써 채비를 마쳤다. 이제 그들에게는 \VUgras{뜨개옷}
\textsuperscript{(2)}과 \VUgras{수영 바지} \textsuperscript{(3)}뿐이다. 그들은 동무들을
기다리며 이야기를 나눈다. 장과, 건너편 기슭에서 벌어지는 잔치 이야기다.

%%K t03-c2-04-1 p
4. --- 그들 곁 땅바닥에는 동무들의 \VUgras{옷}이 놓여 있다. \VUgras{반장화}
\textsuperscript{(4)} 한 \VUgras{켤레}, \VUgras{구두} \textsuperscript{(5)},
\VUgras{양말} \textsuperscript{(6)}, \VUgras{펠트} \textsuperscript{(7)} 모자와
\VUgras{밀짚} \textsuperscript{(8)} 모자, \VUgras{넥타이} \textsuperscript{(9)}다.

%%K t03-c2-05-1 p
5. --- 젊은이 하나는 제 \VUgras{웃옷} \textsuperscript{(10)}을 정성껏 개었다. 그는
그것을 세면 수건 위에 놓는데, 그 곁의 동무도 곧 거기에 자기 \VUgras{옷} 두 벌을 싸
넣을 것이다.

%%K t03-c2-06-1 p
6. --- 여섯째 아이는 늦다. 그는 아직 \VUgras{챙 모자} \textsuperscript{(11)}를 머리에
쓰고 있다. \VUgras{셔츠} \textsuperscript{(12)}도 \VUgras{깃} \textsuperscript{(13)}도
\VUgras{소맷동} \textsuperscript{(14)}도 벗지 않았다. 손에는 제 \VUgras{조끼}
\textsuperscript{(17)}를 들고 있고, \VUgras{바지} \textsuperscript{(16)}와
\VUgras{멜빵} \textsuperscript{(15)}은 아직 그대로다. 다음 경주에 그가 채비를 마치지
못할 것은 뻔하다.

%%K t03-c2-07-1 p
7. --- 젊은이들 곁으로, 양옆에 \VUgras{엉겅퀴} \textsuperscript{(18)}가 많은 오솔길이
강가로 이어진다. 거기에서 야코부스 아저씨가 \VUgras{갈대} \textsuperscript{(19)} 사이에
저녁에 쏘아 올릴 \VUgras{불꽃} \textsuperscript{(20)}을 준비하고 있다.

%%K t03-tit-11 sub
\VUsaut{4.4mm}
\VUpk{10.2pt}{40}{경주.}
\VUsaut{3.4mm}

%%K t03-c2-08-1 p
8. --- 강 위에서는 부인 몇이 양산으로 볕을 가리며 \VUgras{배} \textsuperscript{(21)}로
뱃놀이를 한다. 그들은 아주 볼 만한 경주를 지켜본다. \VUgras{경주 보트}
\textsuperscript{(22)}마다 \VUgras{노 젓는 이} \textsuperscript{(24)} 넷이 온 힘을 다해
젓고, 그동안 \VUgras{키잡이} \textsuperscript{(26)}는 \VUgras{키} \textsuperscript{(25)}
를 잡고 그 여린 배를 몰아간다. \VUgras{노} \textsuperscript{(23)}가 박자에 맞추어 물을
치고, 가벼운 배들은 어지러울 만큼 빠르게 나아간다.

%%K t03-c2-09-1 p
9. --- 젊은이 몇은 다리 기둥 가까이에서 \VUgras{기름 바른 장대} \textsuperscript{(28)}
의 상을 두고 겨룬다. 그 가운데 하나가 휘청거리고 미끄러운 장대 위를 조심스레 걸어간다.
끝에 매달린 작은 \VUgras{깃발} \textsuperscript{(29)}에 닿으려는 것이다. 운 없는
\VUgras{겨루는 이} \textsuperscript{(30)}는 물에 빠졌다. 그는 뭍으로 오려고 헤엄친다.

%%K t03-c2-10-1 p
10. --- 좀 더 나이 든 젊은이들은 \VUgras{배 위 창 겨루기} \textsuperscript{(32)}를
한다 — \VUgras{부두} \textsuperscript{(33)} 가까이에서. 이 모든 놀이를 수많은
\VUgras{구경꾼} \textsuperscript{(31)}이 지켜보는데, 그들은 \VUgras{다리}
\textsuperscript{(27)}의 \VUgras{난간}에 기대거나 \VUgras{강기슭}에 몰려 있다. 그러나
몇몇 구경꾼은 몸을 돌려 \VUgras{상 장대} \textsuperscript{(45)} 놀이를 따라가거나,
상을 \VUgras{나누어 주러} 오는 \VUgras{읍장 행렬}을 바라본다.

%%K t03-c2-11-1 p
11. --- 먼저 \VUgras{개구쟁이} \textsuperscript{(35)} 몇이 있다 — \VUgras{악대}
\textsuperscript{(36)} 앞에 서서 간다. 다음으로 \VUgras{읍장} \textsuperscript{(37)},
그 보좌들, 그리고 작은 읍의 \VUgras{읍 의회}가 온다. 끝으로 \VUgras{소방대원}
\textsuperscript{(38)}들이 걷는다.

%%K t03-tit-12 sub
\VUsaut{5.0mm}
\VUpk{10.2pt}{40}{장.}
\VUsaut{3.6mm}

%%K t03-c2-12-1 p
12. --- \VUgras{장터} \textsuperscript{(39)}에는 사람이 아주 많다. 사람들은 여러
\VUgras{가설 천막}을 보러 온다. 나무 \VUgras{회전목마} \textsuperscript{(40)},
벌써 공연이 시작된 \VUgras{곡마단} \textsuperscript{(41)}, 그리고 읍의 젊은 남녀가
\VUgras{춤}의 즐거움을 누리는 \VUgras{무도장} \textsuperscript{(42)}이다.

%%K t03-c2-13-1 p
13. --- 안쪽으로는 \VUgras{투우장} \textsuperscript{(44)}이 보인다. 거기에서 용감한
스페인 \VUgras{투우사}가 사나운 \VUgras{황소} \textsuperscript{(45)}와 겨룬다. 다음으로
\VUgras{미끄럼 궤도} \textsuperscript{(49)}, \VUgras{사격장} \textsuperscript{(46)}이
있는데, 거기에서는 \VUgras{총} 다루기와 \VUgras{과녁} 맞히기를 익힌다.

%%K t03-c2-14-1 p
14. --- 그 가까이에 \VUgras{장터 극장} \textsuperscript{(47)}이 있어 \VUgras{희극}과
\VUgras{비극}을 올린다. 아주 가까이에는 \VUgras{거인} \textsuperscript{(48)}과
\VUgras{아주 작은 사람}의 천막이 있다. 앞사람의 \VUgras{키}는 이 미터 십오 센티미터요,
뒷사람은 아이만큼도 크지 않다.

%%K t03-c2-15-1 p
15. --- 끝으로 큰 \VUgras{동물 우리} \textsuperscript{(50)}가 있다 — 그 안에는
\VUgras{사나운 짐승}들이 있다. 힘센 \VUgras{발톱}을 가진 \VUgras{호랑이}
\textsuperscript{(51)}, 숱 많은 \VUgras{갈기}를 가진 \VUgras{사자}
\textsuperscript{(52)}, \VUgras{기린} \textsuperscript{(53)} 두 마리, 그리고 제
\VUgras{코}를 그토록 솜씨 좋게 쓰는 \VUgras{코끼리} \textsuperscript{(54)}다.

%%K t03-c2-16-1 p
16. --- 그 짐승들을 길들이는 \VUgras{조련사} \textsuperscript{(55)}가 천막 문 앞에 있다.

\VUsaut{6.0mm}
\VUfilet{22mm}
